\section{Oblivious transfer}\label{ot}

Alice has \(n\) messages \(x_{1},\ldots,x_{n}\). We will assume the
messages are essentially unrelated to each other (since we could always
pad them with random bits). Bob wants to request the \(i\)-th message,
without letting Alice learn anything about the value of \(i\). Alice
wants to send Bob \(x_{i}\), without letting him learn anything about
the other \(n - 1\) messages. An \emph{oblivious transfer (OT)} allows
Alice to transfer a single message to Bob, but she remains oblivious as
to which message she has transferred. We will see two simple protocols
to achieve this.

(In fact, for two-party computation, we only need ``1-of-2 OT'': Alice
has \(x_{1}\) and \(x_{2}\), and she wants to send one of those two to
Bob. But ``1-of-\(n\) OT'' is not any harder, so we will do 1-of-\(n\).)

\subsection{Commutative encryption}

Let us imagine that Alice and Bob have access to some encryption scheme
that is \emph{commutative}:
\[\operatorname{Dec}_{b}( \operatorname{Dec}_{a}( \operatorname{Enc}_{b}( \operatorname{Enc}_{a}(x) ) ) ) = x.\]

In other words, if Alice encrypts a message, and Bob applies a second
layer of encryption to the encrypted message, it does not matter which
order Alice and Bob decrypt the message in -- they will still get the
original message back.

A metaphor for commutative encryption is a box that is locked with two
padlocks. Alice puts a message inside the box, locks it with her lock,
and ships it to Bob. Bob puts his own lock on the box and ships it back
to Alice. What is special about commutative encryption is that Bob's
lock does not block Alice from unlocking her own -- so Alice can remove
her lock and send it back to Bob, and then Bob removes his lock and
recovers the message.

Mathematically, you can get commutative encryption by working in a
finite group (for example \({\mathbb{F}}_{p}^{\times}\), or an elliptic
curve).

\begin{enumerate}
\item
  Alice's secret key is an integer \(a\); she encrypts a message \(g\)
  by raising it to the \(a\)-th power, and she sends Bob \(g^{a}\).
\item
  Bob encrypts again with his own secret key \(b\), and he sends
  \(( g^{a} )^{b} = g^{ab}\) back to Alice.
\item
  Now Alice removes her lock by taking an \(a\)-th root. The result is
  \(g^{b}\), which she sends back to Bob.
\item
  Bob takes a \(b\)-th root, recovering \(g\).
\end{enumerate}

\subsection{OT using commutative encryption}

Our first oblivious transfer protocol is built on the commutative
encryption we just described.

Alice has \(n\) messages \(x_{1},\ldots,x_{n}\), which we may as well
assume are elements of the group \(G\). Alice chooses a secret key
\(a\), encrypts each message, and sends all \(n\) ciphertexts to Bob:
\[\operatorname{Enc}_{a}( x_{1}),\ldots,\operatorname{Enc}_{a}( x_{n} ).\]

But crucially, Alice sends the ciphertexts in order, so Bob knows which
is which.

At this point, Bob cannot read any of the messages, because he does not
know the keys. No problem! Bob just picks out the \(i\)-th ciphertext
\(\operatorname{Enc}_{a}( x_{i} )\), adds his own layer of
encryption onto it, and sends the resulting doubly-encrypted message
back to Alice:
\[\operatorname{Enc}_{b}( \operatorname{Enc}_{a}( x_{i} ) ).\]

Alice does not know Bob's key \(b\), so she cannot learn anything about
the message he encrypted -- even though it originally came from her.
Nonetheless she can apply her own decryption method
\(\operatorname{Dec}_{a}\) to it. Since the encryption scheme is
commutative, the result of Alice's decryption is simply
\[\operatorname{Enc}_{b}( x_{i} ),\] which she sends back to
Bob.

And Bob decrypts the message to learn \(x_{i}\).

\subsection{OT in one step}

The protocol above required one and a half rounds of communication: In
total, Alice sent two messages to Bob (Step 1 and 3), and Bob sent one
message to Alice (Step 2).

We can do better, using public-key cryptography.

Let us start with a simplified protocol that is not quite secure. The
idea is for Bob to send Alice \(n\) keys \[b_{1},\ldots,b_{n}.\]

One of the \(n\), say \(b_{i}\), is a public key for which Bob knows the
private key. The other \(n - 1\) are random garbage.

Alice then uses one key to encrypt each message, and sends back to Bob:
\[\operatorname{Enc}_{b_{1}}( x_{1} ),\ldots,\operatorname{Enc}_{b_{n}}( x_{n} ).\]

Now Bob uses the private key for \(b_{i}\) to decrypt \(x_{i}\), and he
is done.

Is Bob happy with this protocol? Yes. Alice has no way of learning the
value of \(i\), as long as she cannot distinguish a true public key from
a random fake key (which is true of public-key schemes in practice).

But is Alice happy with it? Not so much. A cheating Bob could send \(n\)
different public keys, and Alice has no way to detect it -- like we just
said, Alice cannot tell random garbage from a true public key! And then
Bob would be able to decrypt all \(n\) messages \(x_{1},\ldots,x_{n}\).

But there is a simple trick to fix it. Bob chooses some ``verifiably
random'' value \(r\); to fix ideas, we could agree to use
\(r = \operatorname{sha}(1)\). Then we require that the numbers
\(b_{1},\ldots,b_{n}\) form an arithmetic progression with common
difference \(r\). Bob chooses \(i\), computes a public-private key pair,
and sets \(b_{i}\) equal to that key. Then all the other terms
\(b_{1},\ldots,b_{n}\) are determined by the arithmetic progression
requirement \(b_{j} = b_{i} + (j - i)r\). (Or, if the keys are elements
of a group in multiplicative notation, we could write this as
\(b_{j} = r^{j - i} \cdot b_{i}\).)

Is this secure? If we think of the hash function as a random-number
generator, then all \(n - 1\) ``garbage keys'' are effectively random
values. So now the question is: Can Bob compute a private key for a
given (randomly generated) public key? It is a standard assumption in
public-key cryptography that Bob cannot do this: There is no algorithm
that reads in a public key and spits out the corresponding private key.
(Otherwise, the whole enterprise is doomed.) So Alice is guaranteed that
Bob only knows how to decrypt (at most) one message.

In fact, some public-key cryptosystems (like ElGamal) have a sort of
``homomorphic'' property: If you know the private keys for two different
public keys \(b_{1}\) and \(b_{2}\), then you can compute the private
key for the public key \(b_{2}b_{1}^{- 1}\). (In ElGamal, this is true
because the private key is just the discrete logarithm of the public
key.) So, if Bob could dishonestly decrypt two of Alice's messages, he
could compute the private key for the public key \(r\). But \(r\) is
verifiably random, and it is very hard (we assume) for Bob to find a
private key for a random public key.
